% ==============================================================================
% 1.5 ESTIMACIONES
% ==============================================================================
\subsection{1.5. Estimaciones}

\subsubsection{1.5.1. Dimensiones del proyecto}

\paragraph{1.5.1.1. Tamaño.}

% [PENDIENTE]: Redactar "Tamaño" (26 casos de uso en 6 paquetes; puntos por CU de casos_de_uso_filtrados.md / plan_desarrollo_tareas.md).

\paragraph{1.5.1.2. Complejidad.}

% [PENDIENTE]: Redactar "Complejidad".

\paragraph{1.5.1.3. Estructuración del cliente.}

% [PENDIENTE]: Redactar "Estructuración del cliente".

\subsubsection{1.5.2. Ámbito del proyecto}

\paragraph{1.5.2.1. Objetivos del proyecto.}

Desarrollar un Sistema de Información Web y Móvil para la Postulación y Preselección Automatizada de la Beca Comedor mediante Inteligencia Artificial, con Reportes de Business Intelligence para la Dirección Universitaria de Bienestar Social y Salud (DUBSS) de la Universidad Autónoma Gabriel René Moreno, que permita automatizar la postulación en línea, precalificar algorítmicamente la vulnerabilidad socioeconómica de los estudiantes, gestionar la asistencia física/digital y proveer dashboards analíticos y proyecciones predictivas que apoyen la toma de decisiones institucionales y auditorías.


\paragraph{1.5.2.2. Objetivos específicos.}

\begin{itemize}
    \item Capturar los requisitos funcionales y no funcionales del proceso de postulación y adjudicación de la Beca Comedor mediante entrevistas con el personal de la DUBSS, y especificarlos según el estándar IEEE 830.

    \item Analizar los procesos del sistema modelando los casos de uso y el dominio mediante UML, y definir las reglas de negocio de la preselección.

    \item Diseñar la arquitectura general del sistema web y móvil, y las interfaces de usuario para cada tipo de usuario.

    \item Implementar los módulos del sistema web y móvil, integrando la preselección automatizada con inteligencia artificial y los reportes de Business Intelligence.

    \item Realizar pruebas de caja negra y de caja blanca que verifiquen el cumplimiento de los requisitos especificados.
\end{itemize}


\paragraph{1.5.2.3. Requerimientos principales.}

Los requerimientos del sistema se especifican en las secciones 3.2.1 a 3.2.3 (Especificación de Requisitos de Software según la IEEE 830) del Capítulo 3, y los casos de uso que los implementan, en las secciones 3.2.4 a 3.2.8.

\paragraph{1.5.2.4. Rendimiento.}

% [PENDIENTE]: Redactar "Rendimiento" (RNF: carga <= 3 s, puntaje individual <= 5 s, 20,000 postulaciones por gestión, >= 50 administrativos concurrentes).

\paragraph{1.5.2.5. Fiabilidad.}

% [PENDIENTE]: Redactar "Fiabilidad" (disponibilidad >= 99 % en periodos de postulación y evaluación; idempotencia, RNF-03).

\paragraph{1.5.2.6. Restricciones.}

\begin{itemize}
    \item \textbf{Requisito de Auditoría Física (Restricción Normativa):} El sistema no elimina por completo el uso de papel en la fase de asistencia diaria, debido a la exigencia reglamentaria de los organismos auditores que requieren carpetas físicas con fotocopia de carnet y firma manuscrita para validar la entrega del alimento a cada becado. El software actuará como un medio de control, validación y archivo digitalizador complementario.

    \item \textbf{Operación Interna de la Cocina:} El sistema no gestiona el control de inventarios de insumos, compras de materia prima ni la contabilidad interna de la empresa privada concesionaria del comedor.

    \item \textbf{Comercialización a No Becados:} La venta de almuerzos y platos a estudiantes regulares no becados y personal universitario queda fuera del alcance del sistema, correspondiendo a la gestión propia de la empresa privada.

    \item \textbf{Dependencia de Conectividad:} La carga de expedientes y la sincronización de asistencias requieren conexión a Internet activa.

    \item \textbf{Límites de Verificación Automática:} El sistema realiza una precalificación algorítmica basada en la información declarada y digitalizada; sin embargo, la verificación presencial o visitas domiciliarias mantendrán el criterio operativo del personal de Trabajo Social.
\end{itemize}


\paragraph{1.5.2.7. Infraestructura externa.}

% [PENDIENTE]: Redactar "Infraestructura externa" (Firebase Auth, FCM y Storage; AWS; módulo de IA en Python).

\subsubsection{1.5.3. Estimaciones para el proyecto}

\paragraph{1.5.3.1. Valor esperado.}

% [PENDIENTE]: Redactar "Valor esperado".

\paragraph{1.5.3.2. COCOMO II.}

% [PENDIENTE]: Redactar "COCOMO II".

\paragraph{1.5.3.3. La ecuación del software.}

% [PENDIENTE]: Redactar "La ecuación del software".

\paragraph{1.5.3.4. Planning póker.}

% [PENDIENTE]: Redactar "Planning póker" (sobre casos de uso, no historias de usuario).
